% Source: computational_graph_and _backpropagation.tex
\begin{notechapter}
\chapter{Backpropagation in a Network with Two Hidden Layers}\label{ch:backprop-two-hidden}

\section{Introduction}
This document explains the forward pass, backward pass, and gradient
computations for a neural network with 2 hidden layers. We use row
vectors for inputs and outputs, as this is the convention in most deep
learning frameworks. Additionally, we introduce \textbf{computational
graphs} to visualize the flow of computations and gradients during
backpropagation.

\section{Network Architecture}
Consider a neural network with:
\begin{itemize}[leftmargin=*]
    \item Input layer: 2 neurons (\(x_1, x_2\))
    \item First hidden layer: 3 neurons
    \item Second hidden layer: 4 neurons
    \item Output layer: 1 neuron
\end{itemize}

The weights and biases for each layer are:
\begin{itemize}[leftmargin=*]
    \item First hidden layer: \(\bm{W}_1\) (shape \(2 \times 3\)),
    \(\bm{b}_1\) (shape \(1 \times 3\))
    \item Second hidden layer: \(\bm{W}_2\) (shape \(3 \times 4\)),
    \(\bm{b}_2\) (shape \(1 \times 4\))
    \item Output layer: \(\bm{W}_3\) (shape \(4 \times 1\)),
    \(\bm{b}_3\) (shape \(1 \times 1\))
\end{itemize}

\section{Computational Graphs}
A computational graph is a directed acyclic graph (DAG) where:
\begin{itemize}[leftmargin=*]
    \item \textbf{Nodes} represent variables (inputs, weights, biases,
    intermediate values, outputs).
    \item \textbf{Edges} represent operations (e.g., matrix
    multiplication, addition, activation functions).
\end{itemize}

Computational graphs are used to:
\begin{itemize}[leftmargin=*]
    \item Visualize the flow of computations during the forward pass.
    \item Track dependencies between variables for gradient computation
    during the backward pass.
\end{itemize}

\section{Forward Pass}
The forward pass computes the output of each layer step by step.

\subsection{First Hidden Layer}
\begin{align*}
    \bm{z}_1 &= \bm{x} \cdot \bm{W}_1 + \bm{b}_1 \quad
    \text{(shape } 1 \times 3\text{)} \\
    \bm{a}_1 &= \text{ReLU}(\bm{z}_1) \quad
    \text{(shape } 1 \times 3\text{)}
\end{align*}
Here:
\begin{itemize}[leftmargin=*]
    \item \(\bm{x}\) is the input row vector (\(1 \times 2\)).
    \item \(\bm{W}_1\) is the weight matrix (\(2 \times 3\)).
    \item \(\bm{b}_1\) is the bias row vector (\(1 \times 3\)).
    \item \(\bm{z}_1\) is the pre-activation output (\(1 \times 3\)).
    \item \(\bm{a}_1\) is the activation output (\(1 \times 3\)),
    computed using the ReLU function:
    \[
    \text{ReLU}(\bm{z}) = \max(0, \bm{z})
    \]
\end{itemize}

\subsection{Second Hidden Layer}
\begin{align*}
    \bm{z}_2 &= \bm{a}_1 \cdot \bm{W}_2 + \bm{b}_2 \quad
    \text{(shape } 1 \times 4\text{)} \\
    \bm{a}_2 &= \text{ReLU}(\bm{z}_2) \quad
    \text{(shape } 1 \times 4\text{)}
\end{align*}
Here:
\begin{itemize}[leftmargin=*]
    \item \(\bm{W}_2\) is the weight matrix (\(3 \times 4\)).
    \item \(\bm{b}_2\) is the bias row vector (\(1 \times 4\)).
    \item \(\bm{z}_2\) is the pre-activation output (\(1 \times 4\)).
    \item \(\bm{a}_2\) is the activation output (\(1 \times 4\)),
    computed using the ReLU function.
\end{itemize}

\subsection{Output Layer}
\begin{align*}
    \bm{z}_3 &= \bm{a}_2 \cdot \bm{W}_3 + \bm{b}_3 \quad
    \text{(shape } 1 \times 1\text{)} \\
    \hat{y} &= \sigma(\bm{z}_3) \quad
    \text{(shape } 1 \times 1\text{)}
\end{align*}
Here:
\begin{itemize}[leftmargin=*]
    \item \(\bm{W}_3\) is the weight matrix (\(4 \times 1\)).
    \item \(\bm{b}_3\) is the bias scalar (\(1 \times 1\)).
    \item \(\bm{z}_3\) is the pre-activation output (\(1 \times 1\)).
    \item \(\hat{y}\) is the final output (\(1 \times 1\)), computed
    using the sigmoid function:
    \[
    \sigma(\bm{z}) = \frac{1}{1 + e^{-\bm{z}}}
    \]
\end{itemize}

\section{Loss Computation}
The loss \(L\) is computed using binary cross-entropy loss:
\[
L = -\left( y \log(\hat{y}) + (1 - y) \log(1 - \hat{y}) \right)
\]
where \(y\) is the true label and \(\hat{y}\) is the predicted output.

\section{Backward Pass (Gradient Computation)}
The backward pass computes the gradients of the loss \(L\) with respect
to all weights and biases using the chain rule. Computational graphs
are used to track the flow of gradients.

\subsection{Gradient of the Loss with Respect to the Output Layer}
The gradient \(\frac{\partial L}{\partial \bm{z}_3}\) is computed as:
\begin{align*}
    L &= -\left( y \log(\hat{y}) + (1 - y) \log(1 - \hat{y}) \right) \\
    \hat{y} &= \sigma(\bm{z}_3) = \frac{1}{1 + e^{-\bm{z}_3}} \\
    \frac{\partial L}{\partial \hat{y}} &= -\frac{y}{\hat{y}} +
    \frac{1 - y}{1 - \hat{y}} \\
    \frac{\partial \hat{y}}{\partial \bm{z}_3} &= \hat{y} (1 - \hat{y}) \\
    \frac{\partial L}{\partial \bm{z}_3} &= \frac{\partial L}{\partial \hat{y}}
    \cdot \frac{\partial \hat{y}}{\partial \bm{z}_3} \\
    &= \left( -\frac{y}{\hat{y}} + \frac{1 - y}{1 - \hat{y}} \right)
    \cdot \hat{y} (1 - \hat{y}) \\
    &= \hat{y} - y
\end{align*}
Thus:
\[
\frac{\partial L}{\partial \bm{z}_3} = \hat{y} - y \quad
\text{(shape } 1 \times 1\text{)}
\]

\begin{center}
\begin{tikzpicture}[node distance=2cm, auto]
    \node (L) {\(L\)};
    \node (yhat) [above of=L] {\(\hat{y}\)};
    \node (z3) [above of=yhat] {\(\bm{z}_3\)};
    \node (a2) [left of=z3] {\(\bm{a}_2\)};
    \node (W3) [above of=a2] {\(\bm{W}_3\)};
    \node (b3) [right of=W3] {\(\bm{b}_3\)};

    \draw[->] (yhat) -- (L);
    \draw[->] (z3) -- (yhat);
    \draw[->] (a2) -- (z3);
    \draw[->] (W3) -- (z3);
    \draw[->] (b3) -- (z3);
\end{tikzpicture}
\end{center}

The gradients for \(\bm{W}_3\) and \(\bm{b}_3\) are:
\begin{align*}
    \frac{\partial L}{\partial \bm{W}_3} &= \bm{a}_2^T \cdot
    \frac{\partial L}{\partial \bm{z}_3} \quad
    \text{(shape } 4 \times 1\text{)} \\
    \frac{\partial L}{\partial \bm{b}_3} &= \frac{\partial L}{\partial \bm{z}_3}
    \quad \text{(shape } 1 \times 1\text{)}
\end{align*}

\subsection{Gradient of the Loss with Respect to the Second Hidden Layer}
The gradient \(\frac{\partial L}{\partial \bm{z}_2}\) is computed as:
\begin{align*}
    \frac{\partial L}{\partial \bm{a}_2} &= \frac{\partial L}{\partial \bm{z}_3}
    \cdot \bm{W}_3^T \quad \text{(shape } 1 \times 4\text{)} \\
    \frac{\partial L}{\partial \bm{z}_2} &= \frac{\partial L}{\partial \bm{a}_2}
    \odot \text{ReLU}'(\bm{z}_2) \quad \text{(shape } 1 \times 4\text{)}
\end{align*}
Here, \(\text{ReLU}'(\bm{z}_2)\) is the derivative of the ReLU function:
\[
\text{ReLU}'(\bm{z}_2) = \begin{cases}
    1 & \text{if } \bm{z}_2 > 0 \\
    0 & \text{otherwise}
\end{cases}
\]

\begin{center}
\begin{tikzpicture}[node distance=2cm, auto]
    \node (L) {\(L\)};
    \node (z3) [above of=L] {\(\bm{z}_3\)};
    \node (a2) [above of=z3] {\(\bm{a}_2\)};
    \node (z2) [above of=a2] {\(\bm{z}_2\)};
    \node (a1) [left of=z2] {\(\bm{a}_1\)};
    \node (W2) [above of=a1] {\(\bm{W}_2\)};
    \node (b2) [right of=W2] {\(\bm{b}_2\)};

    \draw[->] (z3) -- (L);
    \draw[->] (a2) -- (z3);
    \draw[->] (z2) -- (a2);
    \draw[->] (a1) -- (z2);
    \draw[->] (W2) -- (z2);
    \draw[->] (b2) -- (z2);
\end{tikzpicture}
\end{center}

The gradients for \(\bm{W}_2\) and \(\bm{b}_2\) are:
\begin{align*}
    \frac{\partial L}{\partial \bm{W}_2} &= \bm{a}_1^T \cdot
    \frac{\partial L}{\partial \bm{z}_2} \quad
    \text{(shape } 3 \times 4\text{)} \\
    \frac{\partial L}{\partial \bm{b}_2} &= \frac{\partial L}{\partial \bm{z}_2}
    \quad \text{(shape } 1 \times 4\text{)}
\end{align*}

\subsection{Gradient of the Loss with Respect to the First Hidden Layer}
The gradient \(\frac{\partial L}{\partial \bm{z}_1}\) is computed as:
\begin{align*}
    \frac{\partial L}{\partial \bm{a}_1} &= \frac{\partial L}{\partial \bm{z}_2}
    \cdot \bm{W}_2^T \quad \text{(shape } 1 \times 3\text{)} \\
    \frac{\partial L}{\partial \bm{z}_1} &= \frac{\partial L}{\partial \bm{a}_1}
    \odot \text{ReLU}'(\bm{z}_1) \quad \text{(shape } 1 \times 3\text{)}
\end{align*}
Here, \(\text{ReLU}'(\bm{z}_1)\) is the derivative of the ReLU function.

\begin{center}
\begin{tikzpicture}[node distance=2cm, auto]
    \node (L) {\(L\)};
    \node (z2) [above of=L] {\(\bm{z}_2\)};
    \node (a1) [above of=z2] {\(\bm{a}_1\)};
    \node (z1) [above of=a1] {\(\bm{z}_1\)};
    \node (x) [left of=z1] {\(\bm{x}\)};
    \node (W1) [above of=x] {\(\bm{W}_1\)};
    \node (b1) [right of=W1] {\(\bm{b}_1\)};

    \draw[->] (z2) -- (L);
    \draw[->] (a1) -- (z2);
    \draw[->] (z1) -- (a1);
    \draw[->] (x) -- (z1);
    \draw[->] (W1) -- (z1);
    \draw[->] (b1) -- (z1);
\end{tikzpicture}
\end{center}

The gradients for \(\bm{W}_1\) and \(\bm{b}_1\) are:
\begin{align*}
    \frac{\partial L}{\partial \bm{W}_1} &= \bm{x}^T \cdot
    \frac{\partial L}{\partial \bm{z}_1} \quad
    \text{(shape } 2 \times 3\text{)} \\
    \frac{\partial L}{\partial \bm{b}_1} &= \frac{\partial L}{\partial \bm{z}_1}
    \quad \text{(shape } 1 \times 3\text{)}
\end{align*}

\section{Weight Updates}
After computing the gradients, update the weights and biases using
gradient descent:
\begin{align*}
    \bm{W}_1 &= \bm{W}_1 - \eta \frac{\partial L}{\partial \bm{W}_1} \\
    \bm{b}_1 &= \bm{b}_1 - \eta \frac{\partial L}{\partial \bm{b}_1} \\
    \bm{W}_2 &= \bm{W}_2 - \eta \frac{\partial L}{\partial \bm{W}_2} \\
    \bm{b}_2 &= \bm{b}_2 - \eta \frac{\partial L}{\partial \bm{b}_2} \\
    \bm{W}_3 &= \bm{W}_3 - \eta \frac{\partial L}{\partial \bm{W}_3} \\
    \bm{b}_3 &= \bm{b}_3 - \eta \frac{\partial L}{\partial \bm{b}_3}
\end{align*}
where \(\eta\) is the learning rate.

\section{Conclusion}
This document provides a step-by-step explanation of forward
propagation, backward propagation, and gradient computations for a
neural network with 2 hidden layers. The use of row vectors simplifies
the notation and aligns with the conventions used in deep learning
frameworks. Computational graphs are introduced to visualize the flow
of computations and gradients, making the process more intuitive.
\end{notechapter}
